\section{Discussion}
\label{sec:discussion}
\subsection{What the result supports}
The primary evidence supports manipulation-conditioned aerial sensing for the
studied retrieval population. The controlled difference is spatial weighting:
the same acquisition model and candidates can direct measurements toward
environment regions that support future manipulator stances rather than toward
generic deficit alone. The observed reduction in no-confirmed-candidate failures
is consistent with this mechanism, while the complete physical endpoint prevents
candidate discovery from being mistaken for successful retrieval.

The method's relevance field should not be read as a success probability. It
combines validated samples, joint-margin quality, footprint geometry and current
operational evidence. It neither models all IK branches nor optimizes navigation,
camera visibility and grasp execution jointly. The one-step gain holds current
support fixed, although future evidence may block or reveal other poses. The
remaining failures, including Generic-only successes, show why a relevance-guided
policy need not dominate on every scene. No guarantee of monotonic task progress
or globally optimal sensing is established.

\subsection{Task, sensing and simulation limitations}
\label{sec:limitations}
The evaluated population contains one known brick, static horizontal ground,
the installed robot geometry and modest perturbations of zero-, one- and
two-wall layouts. Initial sensing begins in a known work region; this is not
an unrestricted aerial search for arbitrary objects. The RGB-D assumptions,
Ground-task reachability map, planar reference bridge and gripper geometry are
task-specific. Uneven terrain, dynamic occlusion, different object geometry and
unseen layouts are outside the present evidence.

Finite-scan opportunity averages a deterministic schedule over starting phases
at a nominal hover pose. It does not calibrate real return likelihood,
reflectance, pose error or intracloud motion distortion. Assumed-height
environment occluders and optimistic unknown-space visibility can mispredict
what the next window actually observes. Environmental evidence also retains
coarse conservative blockers, and finite candidate validation can omit useful
stances. Missing ground support remains a legitimate reason to stop within
budget; predicted opportunities cannot substitute for measurements.

Physics-based execution is stronger evidence than a kinematic reachability
count, but is not real-robot validation. In the current platform, public UAV
map alignment uses a private simulated body-state backend, Ground global
alignment is spawn-derived, and grasp confirmation uses simulated bilateral
contact and joint feedback. The policy receives runtime sensor observations
and public transforms, but the complete localization/contact stack is not
sensor-only in the real-world sense. An independent checker uses simulated
object height to verify lifting. Short retention is additionally supported by
the task's grasp-confirmation hold; the study does not provide an independent
continuous object-height trace throughout a long hold. Dynamics, friction,
calibration and feedback differences limit transfer.

\subsection{Experimental interpretation and transfer}
The 96 paired scenes permit the prespecified primary comparison, not a claim
of universal generalization. Difficulty is a sampled mixture, the qualitative
case is outcome-selected, and auxiliary subsets are small. Conditional
efficiency includes only joint successes, and incomplete distance recordings
preclude a distance-saving claim. The paired success result should therefore
remain the central empirical conclusion, rather than being expanded into
unmeasured energy, speed or real-world robustness benefits.

Future transfer should first establish shared real map/time conventions,
camera/LiDAR extrinsics and hand-eye calibration, then bind the public flight,
base and trajectory interfaces to measured hardware feedback. Real grip and
object-retention evidence require separate characterization; controller
completion alone is insufficient. A stationary sensing/feedback dataset and
offline geometry checks are a suitable first step before bounded motion tests.
No real adapter or hardware experiment is included here. The current simulation
versions and results remain fixed references for any later transfer study.
